\section{Adjoining a marked increment}
\label{mark:section-increments}

Throughout this Section, \(\mathcal C\) is a working scheme whose arrows
are precisely the componentwise strong embeddings preserving its
specified maps and distinguished data.  In particular, there is no
additional condition on intersections or on bases.  We use the working
package of Definition~\ref{lev:working-package}, and the calculus and
isolation results of the preceding two Sections, at sufficiently closed
eligible levels.  All refinements in this Section have fixed set
parameters: the input scheme and template, the prescription below, and,
after it is chosen, the new template.
An unbounded output class is asserted under the eligible-subclass
hypothesis of Convention~\ref{lev:eligible-subclass}.  This supplies both
arbitrarily large closed samples and levels that remain eligible while
having cofinally many smaller template-valid levels.
Theorem~\ref{close:common-levels} supplies the simultaneous closure
bookkeeping.

\begin{definition}[A stationary prescription and its marked pairs]
\label{mark:definition}
A stationary finite prescription for \(\mathcal C\) assigns to every
object \(X\) a type \(p_X\) of a finite tuple, with specified component
sorts.  The assignment commutes with isomorphisms, and, for every arrow
\(X\to X'\), including an arrow between eligible levels, \(p_{X'}\) is
the free extension of \(p_X\), after identifying the base by that arrow.
The prescription is given by fixed set-coded data of the kind used in
the applications.

An object of \(\mathcal C^b\) at level \(\mu\) is a triple
\[
                (X,Y,b),\qquad X\leq Y,
\]
where \(X,Y\) are objects of \(\mathcal C\) at level \(\mu\),
\(b\in Y\) realizes \(p_X\), and \(Y\) is all-dominated over \(X+b\).
The inclusions and the mark are part of the structure.  An arrow is a
commuting pair of componentwise strong embeddings taking the mark to
the mark.  The notation \(X\leq Y\) records an actual inclusion; arbitrary
arrows are put into this form by replacing the codomain with an
isomorphic copy.
\end{definition}

Recall that all-domination quantifies over \emph{every} tuple of length
strictly below the level.  It is not just domination for the members of
one fixed exhaustion of \(Y\).

\begin{lemma}[The pair isomorphism test]
\label{mark:pair-isomorphism}
After a further closure refinement, marked pairs exist over every
object \(X\).  An isomorphism \(e:X\cong X'\) extends to an isomorphism
\[
                 (X,Y,b)\cong(X',Y',b').
\]
More precisely, for bounded tuples \(t\) and \(t'\) in the respective
pair structures, such an extension can take \(t\) to \(t'\) if and only
if \(e\) takes the indicated base entries to their counterparts and
transports
\[
                    \tp(t,b/X)
             \quad\hbox{to}\quad
                    \tp(t',b'/X').
\]
Entries are occurrence-tagged when more than one component is used.

Same-level arrows preserve bounded pair orbits.  Moreover, comparisons
of bounded pair tuples are preserved and reflected when two pairs
at one common small level are mapped to two pairs at a common
larger eligible level.
All tuple lengths and the bounds needed for the comparison are
strictly below the smaller level.
\end{lemma}

\begin{proof}
Existence is Theorem~\ref{iso:all-lengths} applied to the finite
prescription.  The pointed existence and uniqueness assertions follow
from Proposition~\ref{iso:prime}: its back-and-forth form retains prescribed
bounded tuples with the same type together with the mark.  Necessity
of the displayed type condition follows from invariance.

Suppose first that a pair maps to another pair at the same level.
The prescribed type of the mark in the second base is the free
extension of its type in the first base.  For each bounded \(t\) in
the first upper object, Lemma~\ref{iso:free-isolation} therefore says
that its joint type with the mark in the second base is the free
extension of its old joint type.  Choose bounded approximation lists
in the old base computing this joint type and its mark marginal at
the needed testing arities, and include the base entries of \(t\).
An isomorphism of the two bases matching these lists exists by bounded
diagram homogeneity in \(\mathcal C\).  The approximation calculation
and the testing bound show that it transports the whole joint type.
The first part of the Lemma gives an isomorphism of the pairs taking
the old tuple to its image.

The argument also proves preservation when two small pairs with a
bounded match are mapped to a common larger level.  Their joint
types with the marks lift freely by the mixed-level form of
Lemma~\ref{iso:free-isolation}.  Matched small approximation lists
can be matched in the two large bases by underlying homogeneity.
Those lists compute the transported types, and the pair isomorphism
test applies at the large level.

For reflection, let a large pair isomorphism match the two images.
It transports the first small approximation list into the second
large base.  Inside the second small base, realize a copy of this list
over the second small approximation list and the required base
entries.  This is a bounded realization problem: the diagram of the
second list is unchanged under its mixed embedding, and mixed
saturation in the underlying scheme realizes the joint diagram.
Small-level homogeneity extends this match to an isomorphism of the
small bases.

To check that the small isomorphism transports the types, test a
parameter tuple of the size required by the uniform testing bound.
Match that test upstairs over both approximation lists.  Each
list computes the free extension of its small type there, and the
large pair isomorphism identifies the answers.  Transport back to
the small bases.  The answers agree on all tests up to the bound, and
Theorem~\ref{calc:scheme-amalgamation} identifies the whole types.  Apply the small
pair isomorphism test.  All lists used have bounds in the input
arities; closure puts them below the small level.
\end{proof}

\subsection{A template before chain closure}

We first construct a template whose individual values are marked
pairs.  The proof first transfers isolation for tuples with a common
finite bound on the support of each entry, then removes that restriction
by embedding into an all-dominated object.  Neither step uses chain
closure for marked pairs.  That remaining assertion is proved only after
the two-part testing bound is available.

\begin{lemma}[The marked-pair template]
\label{mark:template}
There is a fixed set-sized finitary term template \(F^b\) such that,
on a further class of sufficiently closed eligible levels, its
values on well-orders are objects of \(\mathcal C^b\).  Order
embeddings give the componentwise arrows of Definition~\ref{mark:definition},
also between such levels.  Every marked pair at a level is isomorphic
to a value of this template of that size.

At this stage the following properties are available without using
chain closure in \(\mathcal C^b\): mixed downward
Löwenheim--Skolem, preservation and reflection of bounded diagram
comparisons, bounded diagram saturation, and the exhaustive-list
criterion for being an arrow.
\end{lemma}

\begin{proof}
\par\medskip\noindent\textit{The samples and their language.}\quad
Let \(F\) be the template of \(\mathcal C\).  At arbitrarily high
closed input levels choose a pair
\[
       X=F(J)\leq Y,\qquad Y\cong F(Z),
\]
with \(Y\) all-dominated over \(X+b\).
Choose the presentations sufficiently rich for any prescribed
set-sized order-saturation bound.  We may take \(J,Z\), and a further
order \(W\), to be padded dense orders.  Padding means that the
order has a well-ordered subset of its own cardinality.  In addition,
\(W\) has many disjoint convex substrips, each of the level size.
Proposition~\ref{calc:universal-strip}, in its presentation-invariant
form, puts \(Y\) over \(X\) into
\[
                         L=F(J+W).
\]
All these choices concern the already working scheme \(\mathcal C\).

Expand a sample by sorts for the orders and charts, chart
evaluation maps, the inclusions, and the mark.  For each element of a
chart choose a representing label and its finite list of positions.
Representation tags record the label arity; coordinate functions
have default values outside their arity.  Include a sort of
retraction indices and an evaluation function for a family
\[
                         r_i:Y\longrightarrow X
\]
from Corollary~\ref{calc:retractions}.  Distinct indices are retained as
distinct indices, even if two of their maps happen to be equal.

Finally include finitary LS witness functions for the underlying
\(\K\)-vertices.  For each finite list in a vertex choose a small
strong submodel containing it and the models already chosen for
its proper sublists, and enumerate that model.  Coherence gives
strong inclusions between these small models.  Perform this in all
the finitely many vertex occurrences, including the required
inclusion data.  In the term closure of any generating set, these
enumerated small submodels form a directed cover of each vertex
intersection: every finite tuple has its witness model there,
and a witness for a combined tuple contains the earlier witnesses.
Thus the intersection is strong by the directed-union axioms.
Designated entries keep the necessary sorts nonempty.

This is one fixed set-sized finitary expansion.  It names neither
the diagram table nor all the elements of a high sample.  Its
number of terms has a bound in the old template language and the
fixed input data.

\par\medskip\noindent\textit{Homogeneous finite hulls.}\quad
In each sample take a long ordinal-indexed sequence bundling an
increasing sequence in each order sort and a sequence of distinct
retraction indices.  Color a finite index set by the complete
labelled atomic diagram, with negations and equalities, of its term
hull and the inclusions among the hulls of its subsets.  The palette
has a fixed set bound.  The finite-arity extraction used in
Theorem~\ref{found:order-template} chooses compatible colors at
all finite arities while retaining arbitrarily large homogeneous
samples.  Retain at the same time arbitrarily high prescribed
order-saturation bounds.  Treat the requested homogeneous size
and the requested order saturation as one cardinal demand.
If each possible next color had a bound on the demands it could
meet, the supremum of those set-many bounds would contradict
the availability of samples.  Thus a single next color retains
unbounded demands.  Thinning the horizontal sequence does not
change the saturation of its ambient chart orders.

Let \(H\) be a well-order of a sufficiently closed input size
\(\lambda\), and evaluate the extracted term data on \(H\).
Write the resulting chart orders as \(J_H,W_H,Z_H\), and the
underlying objects as \(X_H,Y_H,L_H\).  The evaluation and
representation functions give both directions of generation:
\[
       X_H=F(J_H),\qquad Y_H\cong F(Z_H),\qquad
       L_H=F(J_H+W_H).
\]
Representation equalities concern fixed finite tuples of terms
and therefore hold in actual homogeneous samples.  The increasing
distinguished sequences pad the chart orders.  Extending the old
template to padded orders uses only the underlying scheme's chain
and directed-union properties.

The vertex LS witnesses prove
\[
                    X_H\longrightarrow Y_H\longrightarrow L_H
\]
componentwise strong over \(X_H\to L_H\).  They also prove this
for maps induced by order embeddings of \(H\), including mixed
sizes.  Each witness model is an isomorphic copy of its sample
model, and the full component map is a directed union of these
strong maps.  Since arrows of \(\mathcal C\) are componentwise,
these are underlying arrows.  All-domination remains to be proved.

\par\medskip\noindent\textit{Bounded horizontal arity.}\quad
Fix a bounded list \(u=(b,c)\) from \(Y_H\) such that each
individual entry has a term representation using at most
\(q\) horizontal indices, for one common finite \(q\).
This is a bound per entry: the union of the supports of
the whole bounded list may be infinite.
Term depths need not be bounded.  Coordinate functions and representations
preserve finite support.  Thus the positions representing \(u\)
in \(J_H+W_H\) use a fixed finite number of indices per position.
The positions representing \(r_i(u)\) in \(J_H\), for the
\(\lambda\) distinguished retraction indices, require only the
additional index \(i\).  Comparing two such positions requires
at most twice a fixed finite arity.

Take an actual sample homogeneous at these arities, with a
horizontal index set containing an order copy of \(H\).
The specified positions match those representing an actual
tuple \(u_*\in Y\) and its images under the corresponding
actual retractions.  Choose the sample with order saturation
above \(\lambda\).  Extend the partial matches to embeddings
of the \emph{entire} orders \(J_H,W_H\) into the respective
sample orders.  In a recursion of at most \(\lambda\) steps,
each next position specifies an order pattern over at most
\(\lambda\) earlier positions, and sample order saturation
realizes it.  The initial partial maps are order embeddings
by the recorded comparisons.  We are not extending a map of
the entire expanded term hull.

The host-order embedding identifies \(u\) with \(u_*\) and
embeds \(X_H\) into \(X\).  The mixed stationary prescription
first shows that the mark in \(Y_H\) realizes \(p_{X_H}\).
We next prove the stronger assertion needed for isolation:
\begin{equation}
       \tp(u_*/X)\text{ is the free extension of }\tp(u/X_H).
       \label{mark:sample-free-descent}
\end{equation}
The images of \(u_*\) under the \(\lambda\) selected actual
retractions belong to the embedded \(X_H\), because their
representing positions were included in the partial matching.
Choose a test bound \(\chi<\lambda\) above the costs for
the arity of \(u\), and retain a set of \(\kappa\) selected
retractions with
\(\chi+\aleph_0<\kappa<\lambda\).
Their images of \(u_*\) form a bounded list in \(X_H\).
For each \(Z\subseteq X\) with \(|Z|\leq\chi\), all but
at most \(|Z|+\aleph_0\) of these retractions fix \(Z\)
pointwise.  At every other index, diagram preservation
identifies the diagram of its image of \(u_*\) with \(Z\)
with the diagram of \(u_*\) with \(Z\).  Thus this image
list computes the actual type over \(X\) with the
majority estimate of Lemma~\ref{calc:average-descent}.
That Lemma gives \eqref{mark:sample-free-descent}.
Only images of the tested tuple have been placed in
\(X_H\); closure of the entire small upper object under
these retractions is neither claimed nor needed here.

The actual sample pair is all-dominated.  Apply
Lemma~\ref{iso:downward-transfer} to
\eqref{mark:sample-free-descent}; its fragment
bound is below \(\lambda\) by closure.  It follows that
\(\tp(b,c/X_H)\) is isolated relative to \(b\).

\par\medskip\noindent\textit{Removing the arity restriction.}\quad
Let \(V\) be all-dominated over \(X_H+b'\), with \(b'\)
of the prescribed type.  We construct an arrow
\[
                       Y_H\longrightarrow V
\]
over \(X_H\), taking \(b\) to \(b'\).
Choose nested bounded lists exhausting \(Y_H\) in countably
many steps.  They may be chosen so that at each step every
entry uses at most a fixed finite number of horizontal indices:
increase a cardinal bound cofinally to \(\lambda\) and the
finite index-arity bound at the same time.  Every term has
finite arity.  Include also full enumerations of small
vertex witness models for finite index sets in the
corresponding bounded part of \(H\).  Their enumerated
entries still have uniformly finite horizontal arity per
entry, independently of term depth.

At each step the whole new list, the preceding list, and
\(b\) satisfy the isolation conclusion already proved.
The constructive primeness argument of
Proposition~\ref{iso:prime} realizes the new list over the
previous assignment and \(X_H\).  Its isolating fragment
together with the preceding list is bounded; underlying
saturation supplies the realization.  Isolation gives
equality of the full types over \(X_H\).  The included
vertex witness models give strongness of the union map.
Every bounded tuple of its image in \(V\) has isolated
joint type with the mark.  Consequently every bounded
tuple of \(Y_H\), with no arity restriction, has this
property.  Thus \(Y_H\) is all-dominated.

\par\medskip\noindent\textit{The remaining pre-closure properties.}\quad
Any marked pair at the level has a base isomorphic to
\(X_H\), and Lemma~\ref{mark:pair-isomorphism} extends this
to a pair isomorphism.  The same Lemma gives orbit
preservation and mixed reflection.  For a set of at most
\(\nu\) entries in a larger pair, take its horizontal
supports and a well-ordered subpresentation of size
\(\nu\) containing them.  Its value is a marked pair by
the individual template assertion just proved; coherence
gives mixed downward Löwenheim--Skolem.

Naturalness on well-orders follows by putting any fixed
bounded patterns into a common well-order enlargement
and using orbit preservation.  Mixed reflection makes
these comparisons independent of the eligible level.
Here are the remaining diagram details, without
assuming the union clause in
Definition~\ref{lev:working-package}.
Define comparison of bounded term patterns by their
pair isomorphism orbits at any eligible level.
The preservation and reflection just proved make
this independent of the level.  At each fixed arity,
the labels and order patterns belong to a fixed
set, giving a uniform orbit bound.  A mixed arrow
factors, by mixed downward Löwenheim--Skolem and
coherence, through a subpresentation of the target
at its source level.  Same-level orbit preservation
in this factor and the level-independent pattern
comparison give mixed diagram preservation.
For saturation, represent a possible joint diagram
in a larger well-ordered pair presentation.
Its support has size below the receiving level,
so that well-ordered support pattern embeds in a
well-order at the receiving level.  Evaluating
\(F^b\) there realizes the joint diagram.
The old subtuple has the orbit of the prescribed
actual old tuple, so an isomorphism of the pairs
carries the new realization over that tuple.
An exhaustive diagram match is componentwise strong by
its small vertex enumerations and respects the
inclusions and mark, so it is a pair arrow.
No union-object assertion for \(\mathcal C^b\) was used.
\end{proof}

\subsection{A testing bound beside an old small pair}

The following form of testing will be used when extending a partial
map.  It separates the use of the pair template from pair-chain
closure, which is still to be proved.

\begin{lemma}[Testing beside an isolated small upper object]
\label{mark:two-part-testing}
Let \(\nu<\mu\) be sufficiently closed levels where
Lemma~\ref{mark:template} holds.  Suppose
\[
       A\leq U,\qquad A\leq C,\qquad C\leq T,T'
\]
are underlying arrows, \(A,U\) are at level \(\nu\), and \(C,T,T'\)
are at level \(\mu\).  Suppose \((A,U,b)\) is a small marked pair
inside \(T\), its specified image \((A,U',b')\) is inside \(T'\),
and the map \(U\cong U'\) fixes \(A\).
Suppose further that \(b/C\) and \(b'/C\) are the prescribed free
extensions and corresponding bounded tuples of \(U,U'\),
jointly with the marks, have equal types over \(C\).

For each \(\sigma<\nu\) there is a bound \(k(\sigma)<\nu\), uniform
in these objects, with the following property.  If \(v\in T\) and
\(v'\in T'\) have length at most \(\sigma\), their comparisons on
joint tests from \(U\) and \(C\) of total size at most
\(k(\sigma)\) determine their comparisons on all such tests of
total size \(<\nu\).  Tests in \(U'\) are read by the specified map.
\end{lemma}

\begin{proof}
Present \((A,U,b)\) as \(F^b(H)\) on a well-order of size \(\nu\).
Given a questioned bounded tuple from \(C\), include it over
\(A\) in a level-\(\nu\) subobject \(V\leq C\).
Append a universal strip to \(H\), retaining size \(\nu\), to
obtain \(H'\).  The base presentation of the pair template is
a natural term presentation into the already working scheme
\(\mathcal C\).  The presentation-invariant universal-strip
result embeds \(V\) over \(A\) into \(A_{H'}\).
Corollary~\ref{calc:mixed-realization} then embeds \(A_{H'}\) into
\(C\) over this copy of \(V\).

The joint bounded diagrams of \(U\) and \(A_{H'}\) inside \(T\)
are the diagrams of the natural layout
\begin{equation}
                         U_H\cup A_{H'}
                 \quad\hbox{inside }U_{H'}.
                 \label{mark:two-part-layout}
\end{equation}
Indeed, each bounded tuple of \(U\), together with \(b\),
is isolated over \(A\), whereas the marginal of \(b\)
over \(A_{H'}\) is the prescribed free extension.
Lemma~\ref{iso:free-isolation} determines its joint type.
The argument applies both to \eqref{mark:two-part-layout}
and to the specified image inside \(T'\).  It includes
equalities between the two parts of the layout; an
additional actual overlap is not ignored.

For a bounded candidate tuple in either ambient, each
bounded test over the layout has a local realization
in \(U_{H'}\), by mixed diagram saturation in
\(\mathcal C\).  Apply Lemma~\ref{calc:compact-placement}
to lift its comparisons to a placed comparison over
\eqref{mark:two-part-layout}.  The layout has two
well-ordered horizontal segments with designated term
parts.  Proposition~\ref{calc:testing} gives a bound depending
on the candidate arity and the fixed template, not on
\(H,H'\) or \(V\).

If the two candidates disagree on the questioned test,
their placed comparisons disagree and yield a witness
within that bound.  Its second-part parameters lie in
\(A_{H'}\), already embedded into \(C\); its first-part
parameters lie in the original \(U\).  This is the
required small witness.  Close \(\nu\) under the
testing and placement costs for this fixed layout.

All padded-order repetitions here concern the
horizontal top presentation as a functor into
\(\mathcal C\).  Its extension to padded orders uses
chain closure of \(\mathcal C\), not of
\(\mathcal C^b\).
\end{proof}

\subsection{Chain closure with all bounded tuple lengths}

\begin{lemma}[Chain closure]
\label{mark:chain-closure}
At further sufficiently closed levels, \(\mathcal C^b\)
has unions of increasing chains of total size equal to
the level.  The union is a marked pair.  The
common-upper-bound clause holds as well, including
into an eligible larger level.
\end{lemma}

\begin{proof}
Fix a level \(\mu\) closed under the bounds of
Lemma~\ref{mark:two-part-testing}, with cofinally many
smaller levels at which Lemma~\ref{mark:template} and
the required underlying results hold.  We do not
require pair-chain closure at the smaller levels.

Write the chain, replacing arrows by inclusions, as
\[
                  (C_i,T_i,b)\quad(i<\alpha),
\]
and its component unions as \(C,T\).  The underlying
scheme gives \(C\leq T\) and its union maps.  The
actual marginal of the finite mark over \(C\) is
\(p_C\): the marginals at proper stages are the
prescribed free extensions, and
Proposition~\ref{calc:local-character} gives stationary
continuity for this finite tuple.

Induct on the chain length.  The shorter-length
assertion allows intermediate limit unions to be
inserted.  Restricting cofinally, we may consider
a continuous chain indexed by a regular cardinal
\(\delta\).  If a bounded tuple \(d\in T\) is
contained in some \(T_i\), its type with \(b\) is
isolated over \(C_i\);
Lemma~\ref{iso:free-isolation} gives isolation
over \(C\).

Suppose \(d\) is not contained in a proper stage.
Its entries occur at cofinally many indices, so
\[
                        \delta\leq |d|<\mu.
\]
Choose an eligible \(\nu<\mu\) strictly above
\(|d|+\delta+\aleph_0\), closed under the arity,
isolation, approximation, and two-part testing
bounds to be used.  Small pair objects below
will come from Lemma~\ref{mark:template} or from
explicit descent; pair-chain closure at \(\nu\)
is not an induction hypothesis.  Choose
\[
             B\geq C
       \quad\hbox{all-dominated over }C+b',
       \qquad b'\models p_C .
\]

It suffices to find an image \(d'\) in \(B\) with
\(\tp(b,d/C)=\tp(b',d'/C)\), because the type on the
right is isolated relative to \(b'\).  We obtain this
image by building small pairs inside the given chain
and compatible maps from their upper objects into
\(B\).  The maps must preserve types over the entire
union base \(C\), including for tuples that span a
limit of the small pairs.

Specifically, we construct continuous systems of underlying
level-\(\nu\) objects \(A_i\leq U_i\) inside
\(C_i\leq T_i\), for \(i<\delta\), and compatible
arrows \(f_i:U_i\to B\), satisfying:
\begin{enumerate}
\item \(b\in U_i\), \(f_i(b)=b'\), and \(f_i\)
fixes \(A_i\);
\item \((A_i,U_i,b)\) is a marked pair;
\item for every tuple \(u\) from \(U_i\) of
length \(<\nu\),
\begin{equation}
                 \tp(b,u/C)
                    =\tp(b',f_i(u)/C);
                 \label{mark:chain-invariant}
\end{equation}
\item every entry of \(d\) is eventually included.
\end{enumerate}
Equation~\eqref{mark:chain-invariant} is computed
in the actual underlying ambients \(T,B\).

\par\medskip\noindent\textit{Retractions prepared for the limits.}\quad
For every nonzero limit \(j\leq\delta\), put
\(C_\delta=C\), and choose in advance families
of \(\nu\) approximate retractions
\begin{equation}
          r^j_\xi:T\to C_j,\qquad
          s^j_\xi:B\to C_j
                    \quad(\xi<\nu).
          \label{mark:limit-retractions}
\end{equation}
They exist in the underlying scheme since
\(C_j\leq T,B\); each old parameter of \(C_j\)
is fixed outside a finite set of indices in
each family.  At a limit we require
\begin{equation}
        r^j_\xi[U_j]\ \cup\
        s^j_\xi[f_j[U_j]]\ \subseteq A_j
                    \quad(\xi<\nu).
        \label{mark:limit-closure}
\end{equation}

We describe the scheduling.  Once an entry
\(x\) and its image \(f_i(x)\) have been
chosen at stage \(i<j\), each value in
\eqref{mark:limit-retractions} belongs to
\(C_j\), hence to some \(C_k\) with \(k<j\)
by continuity.  Assign its inclusion in the
small base to a successor stage \(k'\) with
\(i<k'<j\) and \(k\leq k'\).  Such a stage
exists because \(j\) is a limit.  A value
is never requested before it is available.
New entries and images create at most
\(\nu\cdot\nu\cdot|\delta|=\nu\) requests
at a stage.  The union of the requests
assigned to any one stage still has size
at most \(\nu\): there are \(<\nu\)
stages, each of size \(\nu\).  Underlying
mixed downward Löwenheim--Skolem can
therefore include all of them.  This
also schedules \(j=\delta\), whose
families were chosen in advance.

\par\medskip\noindent\textit{Small objects at a nonlimit stage.}\quad
At an initial or successor stage \(i\),
include the preceding small base
in the new base and the preceding
small top in the new top,
the requests now due, and the designated
entries of \(d\) available in \(T_i\).
Choose \(\nu\) approximate retractions
\(T_i\to C_i\).  Alternate for \(\omega\)
steps the following operations: choose
a level-\(\nu\) subobject of \(C_i\)
containing the current base, all due
base entries, and all the chosen
retraction images of the current top;
then choose a level-\(\nu\) subobject
of \(T_i\) containing that base, the
current top, and required top entries.
Each request has size at most \(\nu\).
Underlying coherence gives increasing
arrows and the base-to-top arrows.
Take underlying unions to obtain
\(A_i\leq U_i\).

The marginal of \(b\) over \(A_i\)
is \(p_{A_i}\), by restriction of the
mixed prescription from \(C_i\).
Lemma~\ref{calc:retraction-descent}
says that every \(<\nu\) tuple from
\(U_i\), jointly with \(b\), has
actual type over \(C_i\) freely
extended from \(A_i\).  Its type
over \(C_i\) is isolated relative
to \(b\), since the given large
pair is all-dominated.  The uniform
fragment bounds are below \(\nu\).
Lemma~\ref{iso:downward-transfer}
proves that \((A_i,U_i,b)\) is a
small marked pair.  Only underlying
unions were used in obtaining it.

\par\medskip\noindent\textit{Extending the map at a nonlimit stage.}\quad
Let \(U^-,A^-,f^-\) be the preceding
small top, base, and map, when present.
The old map fixes every element of
\(U^-\cap C\), not just \(A^-\):
apply \eqref{mark:chain-invariant}
with that element as a parameter.
Thus fixing \(A_i\) is compatible
with the old map on their overlap.

Enumerate \(U_i\) by nested bounded
lists in countably many steps,
including the mark and whole
smaller strong vertex submodels.
Place these lists in \(B\), retaining
the old map, and require equality
on all joint tests of size \(<\nu\)
from \(U^-\) and \(C\).  For a
current source list \(v\), let
\(\chi_0<\nu\) be the bound of
Lemma~\ref{mark:two-part-testing}.
With no preceding stage, use
whole-base testing instead.

For every old test \(u\) from
\(U^-\) of length at most \(\chi_0\),
the tuple \((b,u,v)\) lies in \(T_i\).
Its type over \(C_i\) is isolated,
and its actual mark marginal over
\(C\) is free.  Lemma~\ref{iso:free-isolation}
gives an isolated type over \(C\).
Choose an isolating fragment
\[
              Z_u\subseteq C,\qquad
              |Z_u|\leq h(|b,u,v|)<\nu .
\]
Include the empty old test.  The
union of the fragments has size
at most
\(\nu^{\chi_0}\cdot h(|b|+\chi_0+|v|)\).
This may exceed \(\nu\), but it
is less than \(\mu\), since \(\mu\)
is strong limit and \(\nu<\mu\).

The placement parameters are
\(f^-[U^-]\), the fragments, \(b'\),
previous assignments in the current
exhaustion, and the required fixed
entries of \(A_i\).  Their total
size is \(<\mu\).  Choose a test
size \(\chi<\nu\) large enough
for \(\chi_0\), the current list,
the fragment bounds, all relevant
equalities, and the compact-placement
costs.  Specify joint diagrams on
all parameter tests of size at
most \(\chi\), using the actual
source diagrams.

Every such \emph{simultaneous}
local request is possible in
\(B\).  A query of size at most \(\chi\) may meet many of the
fragments \(Z_u\), but all those entries together form one tuple
from \(C\) of size at most \(\chi\).  Combine that tuple, its old
entries from \(U^-\), fixed base entries, and previous
assignments into one parameter list.  For the first placement,
\eqref{mark:chain-invariant}
matches its old entries over
\(C\); if there is no old stage,
the prescribed mark types give
the initial match.  At later
placement steps, the maintained
joint comparisons give the match
also with previous assigned entries.
The combined list has size \(<\nu\).
Underlying saturation in \(B\)
realizes the diagram of the new
source list over these parameters.

The compact-placement cost depends on the current tuple arity,
\(\chi_0\), the uniform fragment bound, and the fixed labels, rather
than on the size of the entire union of fragments.  For an infinite
\(\sigma<\nu\) bounding the current tuple arity, choose \(\chi<\nu\)
above these costs with \(\chi^\sigma=\chi\).  Each parameter
subset of size at most \(\chi\) then has at most \(\chi\) support
profiles.  The hit-and-separation construction tests only a final
parameter set of size at most \(\chi\).
Apply Lemma~\ref{calc:compact-placement}
to the coherent local requests
in a well-ordered presentation
of \(B\).  It gives a placement
in an order enlargement.  The
whole parameter set has size
\(<\mu\), so underlying saturation
realizes its joint diagram with
that placement back inside \(B\).
In particular the realization
agrees on every \(Z_u\) and all
previously imposed equalities.
Its mark marginal over \(C\)
is \(p_C\).  Isolation gives
the correct whole type over
\(C\) for \((b,u,v)\) when
\(|u|\leq\chi_0\).
Lemma~\ref{mark:two-part-testing}
upgrades this to all joint tests
of size \(<\nu\).  Continue.

The union is an underlying arrow
\(f_i:U_i\to B\), by the included
strong vertex enumerations.
It fixes \(A_i\) and takes \(b\)
to \(b'\).  A \(<\nu\) tuple
of \(U_i\) might span this
countable exhaustion.  For
such a tuple, the completed
arrow gives the correct small
type over \(A_i\).  The small
pair is all-dominated, and
both actual mark marginals
over \(C\) are the prescribed
free extensions.  Mixed
Lemma~\ref{iso:free-isolation}
identifies their types over
\(C\).  This proves
\eqref{mark:chain-invariant}
for every \(<\nu\) tuple.

\par\medskip\noindent\textit{Proper limit stages.}\quad
For a nonzero limit \(j<\delta\),
take underlying unions
\[
 A_j=\bigcup_{i<j}A_i,\qquad
 U_j=\bigcup_{i<j}U_i,\qquad
 f_j=\bigcup_{i<j}f_i .
\]
They are underlying objects
and an underlying arrow by
the chain and mixed upper-bound
clauses of \(\mathcal C\).
Their size is \(\nu\), since
\(|j|<\nu\).  Scheduling gives
\eqref{mark:limit-closure}.
For every \(<\nu\) tuple of
\(U_j\), apply
Lemma~\ref{calc:retraction-descent}
on each side: its actual
type over \(C_j\) is the
free extension of its type
over the common \(A_j\).
These small types agree
because \(f_j\) fixes \(A_j\).

The given large pair
\((C_j,T_j,b)\) is all-dominated.
Downward isolation transfer
proves all-domination over
\(A_j+b\), without small-pair
chain closure.  The equality
over \(C_j\), isolation there,
and prescribed free mark
marginals over \(C\) then
give equality over \(C\) by
Lemma~\ref{iso:free-isolation}.
All proper-stage invariants
are established.

\par\medskip\noindent\textit{The final limit.}\quad
Take underlying unions
\(A_\delta,U_\delta,f_\delta\).
The families prepared for
\(j=\delta\), condition
\eqref{mark:limit-closure},
and Lemma~\ref{calc:retraction-descent}
give equality of the actual
types over \(C\) from their
common type over \(A_\delta\).
No isolation assertion for
the final source pair has
been used.  In particular,
\[
       \tp(b,d/C)
          =\tp(b',f_\delta(d)/C),
\]
since \(|d|<\nu\) and all its
entries were included.
The right-hand type is
isolated by all-domination
of \(B\); hence the left-hand
type is isolated relative
to \(b\).  This covers every
bounded tuple not already
contained in a proper stage,
and proves all-domination
of \(T\) over \(C+b\).

The union maps are componentwise
strong by the underlying
scheme.  If all stages lie
in a common marked pair,
the underlying upper-bound
clause gives a componentwise
strong union map preserving
the data and mark.  That
is precisely a pair arrow.
The argument includes a
larger eligible common upper
bound.
\end{proof}

\begin{theorem}[The marked-increment working package]
\label{mark:package}
Let \(\mathcal C\) be a working scheme
with componentwise arrows and a stationary
finite prescription as in
Definition~\ref{mark:definition}.
Assume the eligible-subclass hypothesis of
Convention~\ref{lev:eligible-subclass}.
On an unbounded further class of sufficiently closed
eligible levels, \(\mathcal C^b\) is again
a working scheme, with componentwise
arrows, fixed template \(F^b\), and mixed
arrows between these levels.  The output class
retains the eligible-subclass hypothesis.

The added language and bound recipes
at this step have set-sized bounds in
the input template and prescription.
Thus the construction can be performed
at each fixed finite depth without
class-sized parameters or a language
size depending on the level.
\end{theorem}

\begin{proof}
Lemma~\ref{mark:pair-isomorphism}
gives the isomorphism test, orbit
preservation, and mixed reflection.
Lemma~\ref{mark:template} gives
individual template values and their
mixed arrows, mixed downward
Löwenheim--Skolem, natural bounded
diagrams and saturation, and the
exhaustive-list arrow criterion.
Lemma~\ref{mark:chain-closure}
supplies the missing union-object
and upper-bound assertions.
Composition and coherence are
componentwise.

The extraction added only finitary
chart, representation, retraction,
and LS-witness symbols with a
fixed set bound.  Further bounds
were testing and placement costs
for the old template and the
fixed two-part layout of the
new template, isolation and
approximation costs, and a
refinement permitting cofinally
many smaller template-valid levels.
These are fixed-parameter
operations in the closure
bookkeeping.  No step requires
small-pair chain closure before
it is proved.
\end{proof}
